\documentclass[11pt]{article}
\usepackage[margin=1in]{geometry}
\usepackage[T1]{fontenc}
\usepackage[utf8]{inputenc}
\usepackage{array}
\usepackage{booktabs}
\usepackage{ragged2e}
\usepackage{tabularx}
\usepackage{hyperref}
\usepackage{xcolor}
\usepackage{listings}
\lstset{basicstyle=\ttfamily\small,breaklines=true,columns=fullflexible,frame=single}
\setlength{\parindent}{0pt}
\setlength{\parskip}{0.65em}
\newcolumntype{L}{>{\RaggedRight\arraybackslash}X}
\newcommand{\profiletableheader}{%
Benchmark & Total insts & Load & Store & \shortstack{Uncond.\\branch} &
\shortstack{Cond.\\branch} & \shortstack{Integer\\compute} &
\shortstack{Floating pt.\\compute} \\
}
\title{Lab 1: ISA Profiling}
\author{Computer Architecture Laboratory}
\date{}
\begin{document}
\maketitle
\section{Introduction}
In this exercise, students run pre-compiled benchmarks to analyze the distribution of instructions they execute. The \texttt{sim-profile} simulator is used to collect instruction class statistics. The goal is to understand the characteristics of different programs and compare the PISA and Alpha ISAs.
All commands in this lab should be run from the cloned repository root.
\begin{lstlisting}
cd /path/to/SimpleScalar
\end{lstlisting}
\section{Part 1: Profiling Alpha Benchmarks}
\subsection{Step 1: Configure the Simulator for Alpha}
Clean the directory:
\begin{lstlisting}
make clean
\end{lstlisting}
Configure for Alpha:
\begin{lstlisting}
make config-alpha
\end{lstlisting}
Build the simulators:
\begin{lstlisting}
make
\end{lstlisting}
\subsection{Step 2: Run the Benchmark Suite}
Execute the following commands. After each run, find \texttt{sim\_num\_insn} and the \texttt{sim\_inst\_class\_prof} section in the output.
\begin{lstlisting}
./sim-profile -iclass benchmarks/anagram.alpha benchmarks/words < benchmarks/anagram.in
\end{lstlisting}
\begin{lstlisting}
./sim-profile -iclass benchmarks/compress95.alpha < benchmarks/compress95.in
\end{lstlisting}
\begin{lstlisting}
./sim-profile -iclass benchmarks/go.alpha 50 9 benchmarks/2stone9.in
\end{lstlisting}
\begin{lstlisting}
./sim-profile -iclass benchmarks/cc1.alpha -O benchmarks/1stmt.i
\end{lstlisting}
In the current repository, \texttt{benchmarks/} is directly under the repository root.
\subsection{Step 3: Data Collection and Analysis}
Record your results in Table 1.
\begin{table}[htbp]
\centering
\caption{Instruction class profile for Alpha benchmark programs}
\label{tab:alpha-benchmarks}
\small
\setlength{\tabcolsep}{3pt}
\renewcommand{\arraystretch}{1.18}
\begin{tabularx}{\textwidth}{@{}L r c c c c c c@{}}
\toprule
\profiletableheader
\midrule
\texttt{anagram.alpha} &  &  &  &  &  &  &  \\
\texttt{go.alpha} &  &  &  &  &  &  &  \\
\texttt{compress95.alpha} &  &  &  &  &  &  &  \\
\texttt{cc1.alpha} &  &  &  &  &  &  &  \\
\bottomrule
\end{tabularx}
\vspace{0.25em}
\footnotesize All instruction-class values except total instructions are percentages.
\end{table}
Answer these questions for each benchmark:
\begin{enumerate}
\item Is the benchmark memory intensive or computation intensive?
\item Is the benchmark mainly using integer computation or floating-point computation?
\item What percentage of instructions are conditional branches?
\item On average, how many instructions execute between conditional branches?
\end{enumerate}
Use:
\begin{lstlisting}
instructions between conditional branches = 100 / conditional branch percentage
\end{lstlisting}
\section{Part 2: Comparing Alpha and PISA ISAs}
This part executes the same smaller test programs on both Alpha and PISA configurations.
\subsection{Section A: Run Alpha Test Programs}
Ensure the simulator is still configured for Alpha. If needed:
\begin{lstlisting}
make clean
make config-alpha
make
\end{lstlisting}
Run:
\begin{lstlisting}
./sim-profile -iclass tests-alpha/bin/test-math
./sim-profile -iclass tests-alpha/bin/test-fmath
./sim-profile -iclass tests-alpha/bin/test-llong
./sim-profile -iclass tests-alpha/bin/test-printf
\end{lstlisting}
Record the results in Table 2.
\begin{table}[htbp]
\centering
\caption{Instruction class profile for Alpha test programs}
\label{tab:alpha-tests}
\small
\setlength{\tabcolsep}{3pt}
\renewcommand{\arraystretch}{1.18}
\begin{tabularx}{\textwidth}{@{}L r c c c c c c@{}}
\toprule
\profiletableheader
\midrule
\texttt{test-math} &  &  &  &  &  &  &  \\
\texttt{test-fmath} &  &  &  &  &  &  &  \\
\texttt{test-llong} &  &  &  &  &  &  &  \\
\texttt{test-printf} &  &  &  &  &  &  &  \\
\bottomrule
\end{tabularx}
\vspace{0.25em}
\footnotesize All instruction-class values except total instructions are percentages.
\end{table}
\subsection{Section B: Run PISA Test Programs}
Reconfigure for PISA:
\begin{lstlisting}
make clean
make config-pisa
make
\end{lstlisting}
Run:
\begin{lstlisting}
./sim-profile -iclass tests-pisa/bin.little/test-math
./sim-profile -iclass tests-pisa/bin.little/test-fmath
./sim-profile -iclass tests-pisa/bin.little/test-llong
./sim-profile -iclass tests-pisa/bin.little/test-printf
\end{lstlisting}
Record the results in Table 3.
\begin{table}[htbp]
\centering
\caption{Instruction class profile for PISA test programs}
\label{tab:pisa-tests}
\small
\setlength{\tabcolsep}{3pt}
\renewcommand{\arraystretch}{1.18}
\begin{tabularx}{\textwidth}{@{}L r c c c c c c@{}}
\toprule
\profiletableheader
\midrule
\texttt{test-math} &  &  &  &  &  &  &  \\
\texttt{test-fmath} &  &  &  &  &  &  &  \\
\texttt{test-llong} &  &  &  &  &  &  &  \\
\texttt{test-printf} &  &  &  &  &  &  &  \\
\bottomrule
\end{tabularx}
\vspace{0.25em}
\footnotesize All instruction-class values except total instructions are percentages.
\end{table}
\section{Final Analysis}
Compare the Alpha and PISA results. Create a histogram using MATLAB, Excel, Python, or another tool comparing a key metric, such as total instruction count for each test program across the two ISAs.
What can you conclude about the two ISAs from your data and plot? For instance, does one ISA accomplish the same task with fewer instructions?
\end{document}
